\section{Ограничения и промежуточные выводы}

По ходу работы я несколько раз менял дальнейшие планы по результатам
коротких проверок. После сравнения с greedy я стал использовать beam 4.
Для модели 002 я снизил LR после воспроизведения срыва обучения.
В поиске на H100 я выбрал BART 8+8/d1024, который в этом бюджете
обучался лучше более глубоких вариантов. На A100 я не увидел выигрыша
от copy-head, зато получил +1.0661 п.п. при замене повторяющихся
данных четырьмя миллионами различных примеров.

После полного обучения выбранного BART я получил 74.2985\% на той же
validation и уложился в ограничение T4: 165.1 секунды на официальный test.
Однако на public моя прежняя полная модель 002 осталась лучше:
735/960 против 729/960 у новой. Поэтому я не считаю более широкий BART
безусловным улучшением своего решения. Рост относительно короткого
Distinct4m не означал роста относительно предыдущей официальной посылки.

При интерпретации результатов я учитываю следующие ограничения:
\begin{itemize}
\item В основных сравнениях использован один seed. Разница одной задачи
между BART и Copy не устанавливает превосходство архитектуры.
\item Ранняя сходимость на 1M не определяет порядок моделей после полного
обучения; архитектуры оценивались в конкретном LR и бюджете.
\item Validation использовалась для выбора моделей, checkpoint и decoding.
Это не независимая итоговая оценка обобщения.
\item Public новой модели оказался ниже public модели 002 на шесть решений.
По одним агрегатам нельзя восстановить парные изменения ответов
или сделать вывод о качестве на пока закрытой private-половине.
\item Разбиение по строкам допускает одинаковые и близкие условия
в разных частях; зависимость примеров не учтена оценкой метрики.
\item Эталонные программы содержат смысловые дефекты. Согласие с их
выводом отличается от правильности относительно условия.
\item Потери на Val-20k и Val-full имеют разные наборы примеров;
их численные значения нельзя использовать как прямое парное сравнение.
\item Ограничение T4 проверено отдельным реальным прогоном финалиста,
а не пересчётом A100 timing или training throughput. Это один замер;
запас около 14.9 секунды не гарантирует то же время в любом runtime.
\end{itemize}

\section{Как я использовал ИИ-агентов}
Я формулировал цели, задавал ограничения и выбирал направления
экспериментов. GPT/Codex я использовал для обсуждения гипотез,
написания и проверки кода, подготовки запусков, анализа журналов
и ведения отчёта. Например, я поручил Codex реализовать общий pipeline,
подготовить сравнения моделей и собрать метрики из сохранённых результатов.
В таблице я перечислил основные задачи агентов и то, как они помогли
мне принимать решения.

\begin{table}[htbp]
\centering\small
\begin{tabularx}{\linewidth}{>{\raggedright\arraybackslash}p{.20\linewidth}YY}
\toprule
Задача агента & Что было сделано & Влияние на решение \\
\midrule
Подготовка данных & Загрузка, очистка комментариев, разбиения, BPE,
контроль labels и padding & Единый воспроизводимый pipeline для моделей \\
Модели и обучение & BART/T5/Copy, DDP, resume, token-normalized loss,
проверки сохранения и генерации & Возможность сопоставлять архитектуры
без переписывания основного pipeline \\
Диагностика и ревью & Воспроизведение срыва 002, проверка LR,
масок cross-attention, compile и AMP & Выбор стабильного LR,
исправление ошибок до интерпретации результатов \\
Эксперименты & Подготовка конфигураций, preflight, мониторинг,
сбор метрик и парных исходов & Быстрый короткий отбор перед полным
обучением; сохранение фактов при сбоях \\
Разбор результатов & Аудит 200+200, контроль эталонов,
таблицы, графики и этот LaTeX-черновик & Уточнение ограничений метрики
и причин выбора следующего эксперимента \\
\bottomrule
\end{tabularx}
\caption{Использование агентов по журналу проекта.}
\end{table}

Я опирался на тесты, короткие реальные прогоны и сохранённые метрики,
которые агент готовил для проверки предложений. Отдельное ревью помогло
найти, например, скрытое включение compile и ошибочную интерпретацию
отсутствующего \code{import math}. При разборе шума я использовал
суждения агента как подсказки для анализа примеров, но не как основание
для создания новых таргетов или изменения знаменателя \pass{}.
Выбор следующего эксперимента и решения о длительном обучении
я оставлял за собой.

\section{Что добавить в следующую версию}
\begin{enumerate}
\item \pending{После открытия private добавить результаты обеих моделей,
если они будут доступны, и уточнить итоговый вывод о качестве.}
\item При необходимости сравнить старую полную модель 002 с новой
полной моделью на одном согласованном отложенном наборе;
разные public и validation метрики не заменяют такое сравнение.
\item Я хочу дополнить текст показательными примерами программ
и сократить вспомогательные детали, чтобы яснее объяснить логику экспериментов.
\end{enumerate}

\clearpage
\appendix
\section{Источники чисел и воспроизведение рисунков}
Журнал работы находится в \code{README.md}, первичные результаты ---
в каталоге \code{reports/}. Сопровождающий \path{report/source_manifest.json}
фиксирует SHA256 и размер JSON-файлов, из которых построены таблицы и рисунки.
Скрипт \path{report/build_figures.py} читает только сохранённые результаты:
он не запускает обучение, инференс или выполнение программ.

Для обновления иллюстраций из корня репозитория используется команда:
\begin{quote}\small\code{python3 report/build_figures.py}\end{quote}
Исходник документа разделён на
файлы в \code{report/sections/}; точка входа --- \code{report/main.tex}.
Сохранённые проверки Stage2 и Stage4 подтверждают идентичность
validation и эталонов, а финальная сводка полного обучения --- exitcode 0
и соответствие экспорта лучшему checkpoint.
Предварительные планы и незавершённые попытки в этом отчёте явно
отделены от фактически завершённых экспериментов.

